
\section*{Заключение}
В ходе выполнения работы решены поставленные задачи: проведён анализ архитектурных ограничений современных систем анализа трафика и существующих библиотек событий, сформулированы функциональные и нефункциональные требования, спроектирована и реализована расширяемая система генерации событий, оптимизированная для встраивания в DPI-конвейеры. Разработанная библиотека сочетает компиляционную типобезопасность, асинхронную lock-free диспетчеризацию и механизмы горячей перезагрузки обработчиков, обеспечивая предсказуемую производительность и низкие задержки в условиях высокой нагрузки.

Научная новизна работы заключается в применении гибридной модели событийного движка, интегрирующей CRTP-идентификацию типов, векторизованную маршрутизацию и NUMA-оптимизированные пулы памяти, что позволяет достичь линейной масштабируемости на многоядерных архитектурах. Практическая значимость определяется возможностью использования системы в качестве унифицированного промежуточного слоя между парсерами трафика и аналитическими модулями, что сокращает время разработки детекторов угроз, упрощает поддержку событийных схем и снижает операционные затраты на интеграцию в существующие SOC-инфраструктуры.

Экспериментальные исследования подтвердили эффективность разработанного решения: пропускная способность достигла 2.14 млн событий/с при задержке p95=82 мкс, что превосходит эталонные аналоги на 28–34\%. Механизм горячей перезагрузки обеспечивает обновление логики за 140 мс без потери данных, а архитектура backpressure гарантирует устойчивость к пиковым нагрузкам. Перспективные направления дальнейших исследований включают интеграцию с фреймворками машинного обучения для автоматической кластеризации аномальных событий, поддержку аппаратного ускорения через eBPF для снижения накладных расходов на уровне ядра ОС, а также разработку стандартизированного формата обмена событиями для межплатформенной совместимости в гетерогенных сетевых инфраструктурах.

\newpage